\documentclass[12pt,a4paper,notitlepage,twocolumn]{article}
\usepackage{amsmath}
\usepackage{amssymb}
\usepackage{amsbsy}
\usepackage{float}
\usepackage[french]{babel}

\usepackage{palatino}

 \usepackage[active]{srcltx}
\usepackage{scrtime}

\newcounter{qnumber}
\setcounter{qnumber}{0}

\newcounter{enumber}
\setcounter{enumber}{0}

\newcommand{\question}{\textbf{(\addtocounter{qnumber}{1}\theqnumber)}\,}
\newcommand{\exercice}{\textbf{\addtocounter{enumber}{1}\textsc{Exercice} \theenumber}.\,}
\setlength{\parindent}{0cm}


\begin{document}

\title{\textsc{Économétrie Approfondie}\\ \small{(Session de rattrapage)}}
\date{12 juin 2026}

\maketitle

\thispagestyle{empty}

\begin{quote}
  \textit{Les réponses non commentées ou insuffisamment détaillées ne seront pas
    considérées. Prenez le temps de faire des phrases.}
\end{quote}

\bigskip


\exercice On suppose que les données sont générées par le modèle suivant~:
\[
  y_i = x_{i,1}\beta_1 + \ldots x_{i,K}\beta_K + \varepsilon_i
\]
où $x_{i,k}$, pour $k=1,\ldots,K$ sont des variables explicatives
déterministes, $\beta_k$, pour $k=1,\ldots,K$, sont des paramètres
réels, $\varepsilon_i$ une variable aléatoire i.i.d. centrée de
variance $\sigma_{\varepsilon}^2$. \question Donner la représentation
matricielle de ce modèle. \question Définir et donner l'expression de
l'estimateur des MCO du vecteur de paramètres $\beta$, que l'on notera
$\hat b$ (en montrant comment on arrive à cette expression et en
explicitant les hypothèses nécessaires pour que cet estimateur
existe). \question Montrer que cet estimateur est sans
biais. \question Calculer la variance de cet estimateur. \question
Énoncer le théorème de Gauss-Markov et montrer que l'estimateur des
MCO est l'estimateur le plus efficace dans la classe des estimateurs
linéaires et sans biais (BLUE).\newline

\setcounter{qnumber}{0}
\bigskip

\exercice On considère le modèle générateur des données~:
\[
y_i = \beta x_i + \varepsilon_i, \quad i = 1, \ldots, N
\]
où $\beta$ est un paramètre réel, $x_i$ une variable explicative
déterministe non nulle, et $\varepsilon_i$ une variable aléatoire
d'espérance nulle. Les erreurs sont indépendantes mais
hétéroscédastiques~:
\[
\mathbb V\left[\varepsilon_i\right] = \sigma^2 w_i
\]
où les $w_i>0$ sont des poids connus et $\sigma^2>0$. Le modèle
empirique est~:
\[
y_i = b\, x_i + \epsilon_i
\]
\question Donner l'expression de l'estimateur des MCO $\hat b$
de $\beta$. Cet estimateur est-il sans biais~? Justifier. \question
Calculer la variance de $\hat b$. \question Comment, en pratique,
pourrait-on détecter la présence d'hétéroscédasticité dans les
données~? Citer au moins un test et en décrire le principe. \question
L'estimateur des MCO est-il efficace~? Justifier. \question Proposer
une transformation des données conduisant à un modèle à erreurs
homoscédastiques, et en déduire un estimateur efficace $\hat
b_{\textsc{mcg}}$ de $\beta$. De quel estimateur s'agit-il~? \question
Calculer la variance de $\hat b_{\textsc{mcg}}$ et la comparer à celle
de $\hat b$. Conclure. \question Que faire si les poids $w_i$ ne sont
pas connus~?\newline

\textbf{\textsc{Rappel}} \textit{(Inégalité de Cauchy-Schwarz) Pour
  tous réels $a_1,\ldots,a_N$ et $c_1,\ldots,c_N$~:
\[
\left( \sum_{i=1}^N a_i c_i \right)^2 \leq \left( \sum_{i=1}^N a_i^2 \right)\left( \sum_{i=1}^N c_i^2 \right)
\]}

\setcounter{qnumber}{0}
\bigskip

\exercice On suppose que le salaire (en logarithme) d'un individu $i$
est généré par le modèle~:
\[
y_i = \beta_0 + \beta_1 x_i + \beta_2 z_i + \varepsilon_i
\]
où $y_i$ est le (log) salaire, $x_i$ le nombre d'années d'études,
$z_i$ une mesure des capacités («~talent~») de l'individu, et
$\varepsilon_i$ une variable aléatoire i.i.d. centrée, de
variance $\sigma_\varepsilon^2$, non corrélée à $x_i$ et $z_i$. Les
variables $x_i$ et $z_i$ sont aléatoires, i.i.d., centrées, de
variances $\sigma_x^2$ et $\sigma_z^2$, et l'on note
$\sigma_{xz}=\mathrm{cov}(x_i,z_i)\neq 0$ (les individus les plus
talentueux font en moyenne des études plus longues). \newline

Malheureusement, l'économètre n'observe pas le talent $z_i$. Il estime
donc le modèle empirique~:
\[
y_i = b_0 + b_1 x_i + u_i
\]
\question Écrire le terme d'erreur $u_i$ du modèle empirique en
fonction de $z_i$ et $\varepsilon_i$. \question Calculer la covariance
entre le régresseur $x_i$ et l'erreur $u_i$. Qu'en déduit-on sur les
propriétés de l'estimateur des MCO~? \question On note $\hat b_1$
l'estimateur des MCO de la pente. Montrer que~:
\[
\text{plim}_{N\rightarrow\infty}\ \hat b_1 = \beta_1 + \beta_2\,\frac{\sigma_{xz}}{\sigma_x^2}
\]
Interpréter et commenter ce résultat (signe du biais, cas où il
disparaît). \question On suppose qu'il existe une variable $g_i$ (par
exemple le nombre d'années d'études du parent de l'individu)
corrélée avec $x_i$ ($\mathrm{cov}(g_i,x_i)=\sigma_{gx}\neq 0$) mais
non corrélée avec $z_i$ ni avec $\varepsilon_i$. Proposer un estimateur
convergent de $\beta_1$ et montrer sa convergence. \question Quelles
sont les deux conditions que doit satisfaire la variable $g_i$ pour
être un bon instrument~? Cet estimateur est-il sans biais~?

\end{document}

%%% Local Variables:
%%% mode: latex
%%% TeX-master: t
%%% End:
